\section{Related Work}

\paragraph{Attacks that deliver adversarial values to a tool.}
One line of work asks how an adversarial value reaches a tool in the first place, treating the tool itself as a black box that faithfully executes whatever it receives. Adversarial suffixes~\cite{zou2023universal} and indirect prompt injection through retrieved content~\cite{greshake2023not} manipulate what the model decides to send; ToolHijacker~\cite{toolhijacker} and MPMA~\cite{wang2026mpma} instead manipulate \emph{which tool} gets selected, by impersonating or out-competing a legitimate tool's description, and WebInject~\cite{wang2025webinject} shows the same idea survives a change of modality by perturbing screenshot pixels rather than text. Les Dissonances~\cite{les-dissonances} pushes delivery downstream further still, letting one tool's output contaminate a sibling tool's arguments within a session. AgentDojo~\cite{debenedetti2024agentdojo} consolidates these attacks and a growing set of defenses into a shared dynamic benchmark. All of this work asks whether an agent's \emph{decision} can be steered; TOFU instead assumes an adversarial value has already arrived at the tool boundary through any such channel (or through mere carelessness, our A1) and asks whether the \emph{implementation} enforces the security property its interface claims -- a defense that stopped every attack above would still leave TOFU's semantic bugs exploitable by a careless caller, and a tool hardened against every bug TOFU finds is still exposed if its description can be spoofed at selection time.

\paragraph{Defenses, and cataloging the {MCP} threat landscape.}
A second line of work defends against exactly these delivery attacks rather than examining what a tool does once invoked: pre-filtering and localizing injected content~\cite{shi2025promptarmor,jia2026promptlocate}, training models to respect a privilege hierarchy over instruction sources~\cite{wallace2024instruction}, and reconstructing least-privilege permission scopes at the orchestration layer~\cite{zhu2025miniscope}. None of this checks that a tool, once invoked through a permitted channel, behaves according to its own documented semantics; Model-Carrying Code~\cite{sekar2003model} is an early, non-LLM precursor to that idea; TOFU's differential oracle plays a similar role but derives the contract automatically from natural-language tool metadata. A parallel body of survey, taxonomy, and infrastructure work catalogs the broader landscape without testing any single deployment directly~\cite{hou2025landscape,li2026towards,trustshiftprobe,aclemcp,web3survey}: MCPXKIT~\cite{guo2026systematic} unifies 31 attack methods into a single framework (IEEE TDSC); Song et al.~\cite{song2025beyond} give the first end-to-end empirical study of the MCP supply chain, from malicious-server upload through popular aggregator platforms to local exploitation; and MCPCorpus~\cite{lin2025mcpcorpus} contributes the data layer underneath all of this, a publicly released snapshot of roughly 14,000 MCP servers with normalized registry and repository metadata. Within that body, Ding et al.~\cite{aclemcp} argue that remote tool use needs attested, expiring capability leases rather than ambient trust, Rostamzadeh et al.~\cite{trustshiftprobe} characterize \emph{staged} trust attacks in which a server behaves benignly until trust is established, and AEGIS~\cite{aegis} targets cross-domain resource abuse specifically. These establish \emph{that} contracts can be violated in many ways; TOFU complements them with a mechanism for deriving and testing a specific tool's contract automatically.

\paragraph{Catalog-bounded scanners and large-scale dynamic measurement.}
The systems closest to TOFU's stated goal each start from a fixed, predetermined list of weaknesses or sinks: static analyses that never execute the target~\cite{mcp-first-glance,mcp-seclint}, a broad regex rule corpus~\cite{agentthreatrules}, and MCP-Scanner~\cite{mcp-scanner}, which does execute tools but only against a catalog of known threats (Cisco AI Defense's unaffiliated, identically-named open-source tool -- which we refer to as Cisco MCP Scanner~\cite{cisco2025mcpscanner} to avoid confusion with the academic tool above -- follows the same catalog-bounded pattern at industrial scale, layering YARA rules, LLM-based description review, and a cloud threat-intelligence API). MCPSecBench~\cite{yang2025mcpsecbench} systematizes this same catalog-driven approach into a benchmark, useful for comparing scanners but bounded, like them, by its enumerated categories. As Table~\ref{tab:motivation} summarizes, none derives its checks from a tool's own argument semantics. The most closely related contemporaneous work instead moves to dynamic, large-scale measurement, and is where TOFU's contribution most needs careful delineation: VIPER-MCP~\cite{viper-mcp} pairs CodeQL taint queries with an LLM-driven agent that evolves natural-language prompts until an attacker-controlled value provably reaches a sink at runtime, reporting 106 zero-days and 67 assigned CVEs across roughly 39{,}884 repositories -- a scale no single-institution study matches. It is the closest system to \project{}, and the distinction is not scale but what a check can express: VIPER-MCP models exactly three taint-style classes (command injection, SSRF, and path traversal) against a fixed inventory of sink APIs, and its authors state that it does not cover logic vulnerabilities, authentication bypasses, or information disclosure, attributing residual false negatives to unmodeled classes. A taint model is bounded by its sink list; \project{}'s checks are bounded instead by what a tool documents about its own arguments, which is why properties with no corresponding sink -- a destination restricted to an allow-list, a path confined to a declared workspace, an argument whose permitted values the description enumerates -- fall inside \project{}'s reach and outside a taint formulation. Chen et al.~\cite{mcpzoo} study runtime MCP servers at scale and show that existing scanners' verdicts disagree far more than any individual paper suggests -- a warning that applies to \project{} too. Toeppe et al.~\cite{toeppe2026largescale} mine vulnerability patterns from repository history at similar scale, and Padilla~\cite{padilla2026exposed} dynamically assesses the attack surface that internet-facing MCP servers expose by default. A further group intervenes at or after execution rather than before it: FlowGuard~\cite{flowguard} turns runtime signals into evidence for detection, He et al.~\cite{mtguard} combine static and dynamic analysis to secure tool use, MCP-Guard~\cite{xing2026mcpguard} layers a multi-stage defense-in-depth proxy in front of LLM--tool interactions, and Shi et al.~\cite{spellsmith} approach the same taint-style flaws from the opposite direction, rewriting tool \emph{descriptions} to be security-aware so the agent is steered away from unsafe invocations in the first place. What unites this cluster is that its checks generalize across tools precisely because they do not model any individual tool's documented semantics -- the same property that buys them throughput also fixes, in advance, the set of weaknesses they can express. \project{} takes the opposite trade: each of the forty-five sanitizers in Appendix~C is instantiated per argument from that tool's own contract, which requires the LLM-assisted mapping step of Section~4 and does not scale to tens of thousands of repositories. We report this as a difference in what is optimized rather than a claim of superiority, and note that the two are complementary in practice: a taint-style scanner and a semantics-derived one fail on disjoint bugs.

\paragraph{Description-code inconsistency, fuzzing, and LLM-assisted testing.}
A more recent thread asks a narrower version of \project{}'s question directly: Li et al.~\cite{li2026don} show empirically that MCP agents misbehave when tool descriptions are misleading, and Shi et al.~\cite{shi2026description} measure and detect description-code inconsistency in real-world MCP servers, drawing out its security implications. Both substantiate our A4 as common rather than hypothetical, but target description-code mismatch specifically; TOFU sits downstream, asking whether an implementation upholds a documented \emph{security} property regardless of whether it otherwise matches its description. On the mechanical side, TOFU builds on conventional fuzzing infrastructure, using Atheris~\cite{atheris} to mutate structured, schema-valid invocations -- but schema validity alone does not establish safety, and classical crash-based oracles miss exactly the class of bug TOFU targets: a tool that terminates normally, produces a plausible result, and still violates a documented property, as our case study in Section~\ref{sec:casestudy} illustrates. Separately, LLM-assisted fuzzers such as TitanFuzz~\cite{deng2023large} and Fuzz4All~\cite{xia2024fuzz4all} use an LLM to test conventional software; TOFU inverts the relationship (its targets are LLM tools) and narrows the model's role deliberately, restricting it to sanitizer selection and seed generation rather than vulnerability judgment, so the oracle itself stays grounded in an explicit, executable specification.
